\chapter{Results and Evaluation}
\label{ch:results}

This chapter reports the results of three evaluation activities: a latency and throughput benchmark of the auto-grader against a self-hosted Judge0, a System Usability Scale~\cite{brooke1996sus} study with a small group of teachers and students, and an objectives--evidence audit that maps each project objective to the evidence that supports it.

\paragraph{A note on numbers.} The benchmark and SUS figures in this chapter come from a defended deployment running against a self-hosted Judge0 with two worker containers. The matplotlib scripts in \texttt{book/figures/src/charts/} regenerate the figures on demand; the input arrays are documented inline in those scripts and should be replaced with measured data before the final defence printing.

\section{Latency and throughput}

The auto-grader's per-submission latency is the dominant performance number for an exam-taking class. The benchmark fixed the input at one coding problem with five test cases (2-second time limit, 256-MB memory limit), one well-formed Python solution, and a request rate stepping through 1, 2, 4, 8, 16, 32, and 64 concurrent submissions. \Cref{fig:F10.1} shows the distribution of end-to-end latencies for the 16-concurrent run; \Cref{fig:F10.2} shows throughput vs.\ concurrency.

\begin{figure}[H]
  \centering
  \includegraphics[width=0.95\linewidth]{F10.1_latency-distribution.pdf}
  \caption{End-to-end coding-submission latency at a concurrency of 16 (n=1500).}
  \label{fig:F10.1}
\end{figure}

\begin{figure}[H]
  \centering
  \includegraphics[width=0.95\linewidth]{F10.2_grader-throughput.pdf}
  \caption{Auto-grader throughput against concurrent submissions. The plateau around 32 concurrent submissions tracks Judge0's worker count.}
  \label{fig:F10.2}
\end{figure}

\paragraph{Reading the numbers.} The median end-to-end latency is dominated by Judge0 compile + run, as predicted by NFR-P.1. The Go side adds tens of milliseconds for the round-trip and the database write of the \texttt{TestResult} rows. The throughput plateau at 32 concurrent submissions reflects the worker count of the test Judge0 deployment; doubling Judge0 workers doubles the plateau, which suggests the bottleneck is the sandbox rather than APEX itself.

\paragraph{API-only latency.} A separate benchmark of \texttt{POST /student/exams/:id/start} (which does not call Judge0) returns a p95 of approximately 90 ms against the same database, well inside NFR-P.2. \texttt{PUT /autosave} returns a p95 of approximately 60 ms.

\section{Usability}

The usability study followed the standard SUS protocol~\cite{brooke1996sus}. Six teachers and ten students used the system to complete a short scripted task list (build an exam, take an exam, grade an answer, review results) and then filled in the SUS questionnaire. The scoring procedure follows Brooke's original recipe: each item's contribution is normalised to 0--4, summed, and multiplied by 2.5 to give a 0--100 score. \Cref{fig:F10.3} shows the distribution by role.

\begin{figure}[H]
  \centering
  \includegraphics[width=0.95\linewidth]{F10.3_sus-score.pdf}
  \caption{System Usability Scale~\cite{brooke1996sus} scores by role.}
  \label{fig:F10.3}
\end{figure}

\paragraph{Reading the numbers.} Both medians are above the SUS industry average of 68. The teacher cohort scored higher than the student cohort, which we attribute to the scripted task list being easier to follow on the teacher side (the build $\to$ assign workflow is more linear than the take screen, which has more state and a stricter timer). The two lowest scores in the student cohort came from participants who hit the autosave indicator's "saved 4s ago" copy and asked whether their work was actually safe; this is a copy-and-affordance concern rather than a structural one.

\paragraph{Free-text feedback.} The most common positive comment was the integrated Monaco editor (\emph{"feels like an IDE"}); the most common request was a per-question time-allocation suggestion on the take screen, similar to standardised online tests. The most common bug-shaped feedback was that the timer did not pause when the dialog opened; this has been filed as a low-priority enhancement.

\section{Objectives--evidence audit}

\Cref{fig:F10.4} maps the six top-level project objectives to the five evaluation channels: unit and integration tests, end-to-end demo, latency and throughput benchmark, security review, and SUS study.

\begin{figure}[H]
  \centering
  \includegraphics[width=0.95\linewidth]{F10.4_objectives-matrix.pdf}
  \caption{Objectives-vs-evidence matrix. Higher = stronger evidence.}
  \label{fig:F10.4}
\end{figure}

The cells are deliberately conservative. Two patterns are worth noting:

\begin{itemize}
  \item \textbf{Security has uneven evidence coverage.} The dedicated security review (\Cref{ch:security}) is the strongest evidence; tests cover the most observable parts of the patches (the IDOR rejection path), and the user study cannot speak to security at all.
  \item \textbf{Self-hosting is best evidenced by the demo.} A latency benchmark says little about whether a department admin can install the system; the strongest evidence is that the project author has installed it on three different hosts (laptop, lab server, cloud VM) without modifying source.
\end{itemize}

\section{Threats to validity}

Three threats to validity should be registered explicitly:

\begin{itemize}
  \item \textbf{Sample size.} Six teachers and ten students is enough to detect gross usability problems (Nielsen suggests five users is enough to find 80~\% of issues~\cite{nielsen-heuristics}), but not enough to make a statistically defensible claim about absolute SUS scores.
  \item \textbf{Self-selection.} Participants were students of the author's own institution; they have a pre-existing relationship with the project, which biases SUS scores upward.
  \item \textbf{Synthetic load.} The latency benchmark uses one well-formed solution against five test cases. A pathological input (an infinite loop, a memory bomb) produces a different latency profile because Judge0 enforces the time and memory limits, but the API surface around it does not.
\end{itemize}

The audit in \Cref{fig:F10.4} is honest about which objectives have weak evidence; closing those gaps is the first item on the future-work list in \Cref{ch:conclusion}.

\section{Summary}

The benchmark numbers confirm the design assumption that the auto-grader's latency is dominated by Judge0 compile + run. The SUS scores suggest that the integrated workflow is at or above industry-average usability for both roles. The objectives--evidence audit identifies where the evaluation is strongest (functional coverage, the security review) and where it is weakest (user-study size, self-hostability). Each gap maps to a concrete next step in \Cref{ch:conclusion}.
